\chapter{Aprendizado de Máquina Supervisionado}

\section{Por que aprendizado de máquina em uma linguagem de econometria?}

A econometria clássica é forte quando o objetivo é inferência:
estimar um parâmetro, testar uma hipótese ou recuperar um efeito
causal. O aprendizado de máquina é mais útil quando o objetivo é
previsão: encontrar uma função flexível $\hat f(X)$ que aproxime $Y$
bem fora da amostra.

\hayashi{} faz a ponte entre os dois. Seus comandos de ML usam a mesma
sintaxe de fórmula do OLS ou logit e retornam objetos compatíveis com
`predict`, `glance` e `tidy`.

Este capítulo cobre métodos baseados em árvores, redes neurais e um
processo gaussiano. Não é um curso completo de ML; mostra como estimar,
interpretar e comparar esses modelos dentro do \hayashi{}.

\section{Random Forest}

Random Forest (Breiman, 2001) ajusta um conjunto de árvores de decisão
em amostras bootstrap, média suas previsões para reduzir variância.

\begin{lstlisting}[caption={Regressão com Random Forest}]
set_seed(42)
let n = 500
generate df x1 = rnormal(n)
generate df x2 = rnormal(n)
generate df y = 2 + 0.5*x1 - 0.3*x2 + 0.4*x1*x2 + rnormal(n)

let m_rf = rf(y ~ x1 + x2, df, trees=200, depth=10)
print(m_rf)

let g = glance(m_rf)
print(g)
\end{lstlisting}

Principais saídas:

\begin{itemize}
\item \texttt{r2} e \texttt{oob\_r2}: $R^2$ dentro da amostra e out-of-bag
\item importância de variáveis: redução de impureza atribuível a cada
      regressor
\end{itemize}

Use `predict(df, yhat = m_rf)` para obter valores ajustados em um
DataFrame.

\section{Gradient Boosting e XGBoost}

Gradient Boosting constrói árvores sequencialmente, cada uma ajustando
os resíduos do conjunto anterior. XGBoost (Chen e Guestrin, 2016)
extende isso com informação de segunda ordem e regularização.

\begin{lstlisting}[caption={Regressão com XGBoost}]
let m_xgb = xgboost(y ~ x1 + x2, df, trees=100, depth=4, lr=0.1)
print(glance(m_xgb))
\end{lstlisting}

Boosting tende a superar Random Forest em dados tabulares estruturados,
mas é mais sensível a hiperparâmetros e pode sofrer overfit.

\section{Perceptron Multicamadas}

O comando `mlp` ajusta uma rede feed-forward com uma camada oculta,
função de ativação sigmoide e backpropagation com momentum.

\begin{lstlisting}[caption={Rede neural}]
let m_mlp = mlp(y ~ x1 + x2, df, hidden=20, lr=0.01, epochs=500)
print(glance(m_mlp))
\end{lstlisting}

É útil para capturar interações não lineares, mas oferece menos
transparência que árvores e geralmente precisa de ajustes.

\section{LSTM e Transformer para séries}

Para séries temporais, `lstm` e `transformer` usam janelas de
observações para prever o próximo valor. Eles padronizam o alvo
internamente.

\begin{lstlisting}[caption={Previsão com LSTM}]
let n = 200
generate df y = 0.1*(_n - n/2) + 0.5*sin(2*pi()*_n/20) + rnormal(n)
let m_lstm = lstm(df, y, seqlen=10, hidden=10, epochs=100, forecast=5)
print(m_lstm)
\end{lstlisting}

Esses modelos são determinísticos uma vez fixada a semente, mas sua
qualidade preditiva é sensível a convenções. São experimentais.

\section{Processo Gaussiano}

Um Processo Gaussiano define uma distribuição sobre funções. É não
paramétrico no sentido de que não fixa uma forma funcional e reporta
incerteza em torno de cada previsão.

\begin{lstlisting}[caption={Regressão com Processo Gaussiano}]
let m_gp = gp(y ~ x1 + x2, df)
print(m_gp)
\end{lstlisting}

GP é atraente para conjuntos pequenos e quando a propagação de
incerteza importa. É mais lento e mais intensivo em memória que métodos
baseados em árvores.

\section{Quando preferir qual?}

\begin{itemize}
\item \textbf{Interpretabilidade + robustez:} Random Forest ou Gradient Boosting
\item \textbf{Acurácia preditiva com ajuste:} XGBoost
\item \textbf{Funções não lineares suaves:} MLP ou GP
\item \textbf{Séries:} LSTM / Transformer
\item \textbf{Incerteza:} GP ou BART
\end{itemize}

Todos os comandos aceitam uma fórmula padrão do \hayashi{} e retornam um
objeto compatível com `predict` e `glance`. Erros padrão geralmente não
são reportados porque esses métodos não são projetados para inferência
direta sobre coeficientes individuais.
